% GENERATED FILE. Edit canonical lessons, then run npm run generate:books.
% canonical-source-hash: fnv1a64:6add1c2237da2ef4
% canonical-lessons: RU-C174-sovetovat, RU-C174-opisyvat, RU-C174-povtoryat, RU-C174-perevodit, RU-C174-pomnit

\chapter{To advise, To describe, To repeat, To translate, To remember}
\label{ch:ru-to-advise-to-describe-to-repeat-to-translate-to-remember}
\hlchaptermodality{174}

\begin{chapteropening}
\textbf{By the end of this chapter:} \emph{I can say to advise, to describe, to repeat, to translate and to remember.}

Complete the last lesson of chapter 174: I can say to advise, to describe, to repeat, to translate and to remember.
\end{chapteropening}

\section[\texorpdfstring{\ru{советовать} (sovétovat)}{советовать (sovétovat)}]{\ru{советовать} (sovétovat) — to advise}

\label{lesson:RU-C174-sovetovat}



\begin{quote}
Before the new one: say the Russian for to fry, then the Russian for to pour.
\end{quote}

\subsection*{You'll want to know: \ru{советовать}}
\textbf{\ru{советовать}} (\emph{sovétovat}) — \textquotedblleft{}to advise\textquotedblright{}.

\begin{grammarlens}[title={What you've built}]
One more word for this chapter.
\end{grammarlens}

\subsection*{Your turn}
\begin{itemize}
\raggedright
  \item \emph{Say it:} \emph{sovétovat}
  \item \emph{Say it:} \emph{sovétovat}, once more
  \item \emph{Say it:} say \emph{sovétovat}
  \item \emph{Recall:} say the Russian for to fry, then the Russian for to pour, then say \emph{sovétovat} again
\end{itemize}

\begin{tcolorbox}[breakable,colback=teal!4,colframe=teal!35!black,title={Before you move on}]
What does \ru{советовать} mean? (\textquotedblleft{}To advise\textquotedblright{}.) Say it once more.
\end{tcolorbox}
\section[\texorpdfstring{\ru{описывать} (opísyvat)}{описывать (opísyvat)}]{\ru{описывать} (opísyvat) — to describe}

\label{lesson:RU-C174-opisyvat}



\begin{quote}
Before the new one: say the Russian for to pour, then the Russian for to advise.
\end{quote}

\subsection*{You'll want to know: \ru{описывать}}
\textbf{\ru{описывать}} (\emph{opísyvat}) — \textquotedblleft{}to describe\textquotedblright{}.

\begin{grammarlens}[title={What you've built}]
One more word for this chapter.
\end{grammarlens}

\subsection*{Your turn}
\begin{itemize}
\raggedright
  \item \emph{Say it:} \emph{opísyvat}
  \item \emph{Say it:} \emph{opísyvat}, once more
  \item \emph{Say it:} say \emph{opísyvat}
  \item \emph{Recall:} say the Russian for to pour, then the Russian for to advise, then say \emph{opísyvat} again
\end{itemize}

\begin{tcolorbox}[breakable,colback=teal!4,colframe=teal!35!black,title={Before you move on}]
What does \ru{описывать} mean? (\textquotedblleft{}To describe\textquotedblright{}.) Say it once more.
\end{tcolorbox}
\section[\texorpdfstring{\ru{повторять} (povtoryát)}{повторять (povtoryát)}]{\ru{повторять} (povtoryát) — to repeat}

\label{lesson:RU-C174-povtoryat}



\begin{quote}
Before the new one: say the Russian for to advise, then the Russian for to describe.
\end{quote}

\subsection*{You'll want to know: \ru{повторять}}
\textbf{\ru{повторять}} (\emph{povtoryát}) — \textquotedblleft{}to repeat\textquotedblright{}.

\begin{grammarlens}[title={What you've built}]
One more word for this chapter.
\end{grammarlens}

\subsection*{Your turn}
\begin{itemize}
\raggedright
  \item \emph{Say it:} \emph{povtoryát}
  \item \emph{Say it:} \emph{povtoryát}, once more
  \item \emph{Say it:} say \emph{povtoryát}
  \item \emph{Recall:} say the Russian for to advise, then the Russian for to describe, then say \emph{povtoryát} again
\end{itemize}

\begin{tcolorbox}[breakable,colback=teal!4,colframe=teal!35!black,title={Before you move on}]
What does \ru{повторять} mean? (\textquotedblleft{}To repeat\textquotedblright{}.) Say it once more.
\end{tcolorbox}
\section[\texorpdfstring{\ru{переводить} (perevodít)}{переводить (perevodít)}]{\ru{переводить} (perevodít) — to translate}

\label{lesson:RU-C174-perevodit}



\begin{quote}
Before the new one: say the Russian for to describe, then the Russian for to repeat.
\end{quote}

\subsection*{You'll want to know: \ru{переводить}}
\textbf{\ru{переводить}} (\emph{perevodít}) — \textquotedblleft{}to translate\textquotedblright{}.

\begin{grammarlens}[title={What you've built}]
One more word for this chapter.
\end{grammarlens}

\subsection*{Your turn}
\begin{itemize}
\raggedright
  \item \emph{Say it:} \emph{perevodít}
  \item \emph{Say it:} \emph{perevodít}, once more
  \item \emph{Say it:} say \emph{perevodít}
  \item \emph{Recall:} say the Russian for to describe, then the Russian for to repeat, then say \emph{perevodít} again
\end{itemize}

\begin{tcolorbox}[breakable,colback=teal!4,colframe=teal!35!black,title={Before you move on}]
What does \ru{переводить} mean? (\textquotedblleft{}To translate\textquotedblright{}.) Say it once more.
\end{tcolorbox}
\section[\texorpdfstring{\ru{помнить} (pómnit)}{помнить (pómnit)}]{\ru{помнить} (pómnit) — to remember}

\label{lesson:RU-C174-pomnit}



\begin{quote}
Before the new one: say the Russian for to repeat, then the Russian for to translate.
\end{quote}

\subsection*{You'll want to know: \ru{помнить}}
\textbf{\ru{помнить}} (\emph{pómnit}) — \textquotedblleft{}to remember\textquotedblright{}.

\begin{grammarlens}[title={What you've built}]
One more word for this chapter.
\end{grammarlens}

\subsection*{Your turn}
\begin{itemize}
\raggedright
  \item \emph{Say it:} \emph{pómnit}
  \item \emph{Say it:} \emph{pómnit}, once more
  \item \emph{Say it:} say \emph{pómnit}
  \item \emph{Recall:} say the Russian for to repeat, then the Russian for to translate, then say \emph{pómnit} again
\end{itemize}

\begin{tcolorbox}[breakable,colback=teal!4,colframe=teal!35!black,title={Before you move on}]
What does \ru{помнить} mean? (\textquotedblleft{}To remember\textquotedblright{}.) Say it once more.
\end{tcolorbox}
